\begin{lemma}
Let \(I\) be a \(k\)-set and assume the degree event
\(\noderef{FiberAndDegreeEvent}\), the medium closed-pairs event, and the small
closed-pairs event.  Suppose also that the large contribution
\[
|\noderef{TwoBiteClosedPairsInClass}(I,L_I)|
\]
is \(\varepsilon_1\)-small in the sense of
\(\noderef{ClosedPairsBound}\), and that the sum of the same-color huge
projected contributions
\[
\sum_{x\in H_I\cap V_R}\binom{|\pi_R(X_x(I))|}{2}_{\mathbb R}
+
\sum_{x\in H_I\cap V_B}\binom{|\pi_B(X_x(I))|}{2}_{\mathbb R}
\]
is also \(\varepsilon_1\)-small.  Then the total contribution from
\[
L_I,\qquad M_I,\qquad S_I,\qquad
\text{and the same-color projections of }H_I
\]
is at most \(4\varepsilon_1k^2\), equivalently it satisfies
\(\noderef{ClosedPairsBound}\) with parameter \(4\varepsilon_1\) and scale
\(k\).
\end{lemma}
\begin{proof}
Fix the data appearing in the statement:
\(\omega\), \(I\), \(\varepsilon,\varepsilon_1,\varepsilon_2\), and the
integer \(k\) with \(|I|=k\).  Define the four real quantities, in exactly the
order in which they occur in the Lean conclusion, by
\[
A_L=|\noderef{TwoBiteClosedPairsInClass}(I,L_I)|,\quad
A_M=|\noderef{TwoBiteClosedPairsInClass}(I,M_I)|,\quad
A_S=|\noderef{TwoBiteClosedPairsInClass}(I,S_I)|
\]
and
\[
\begin{aligned}
H_R(x)&\Longleftrightarrow
  \noderef{TwoBiteHugeClass}(\omega,I,x)\text{ and }x\in V_R,\\
H_B(x)&\Longleftrightarrow
  \noderef{TwoBiteHugeClass}(\omega,I,x)\text{ and }x\in V_B,
\end{aligned}
\]
so that \(H_R\) and \(H_B\) are the natural-language versions of the
let-bound predicates
\[
\begin{aligned}
\mathsf{hugeRed}(x)&:=
  \noderef{TwoBiteHugeClass}(\omega,I,x)\wedge
  \exists r,\ x=\operatorname{red}(r),\\
\mathsf{hugeBlue}(x)&:=
  \noderef{TwoBiteHugeClass}(\omega,I,x)\wedge
  \exists b,\ x=\operatorname{blue}(b)
\end{aligned}
\]
in the Lean statement.  Put
\[
A_H^{=}
=
\noderef{TwoBiteProjectedPairSum}(\omega,I,\mathsf{hugeRed},\pi_R)
+
\noderef{TwoBiteProjectedPairSum}(\omega,I,\mathsf{hugeBlue},\pi_B).
\]
This is exactly the final parenthesized let-bound term in the Lean
conclusion: the first summand ranges over huge red base vertices and uses
the red projection \(\pi_R\), and the second ranges over huge blue base
vertices and uses the blue projection \(\pi_B\).

Unfolding \(\noderef{ClosedPairsBound}\), the large-class hypothesis is
precisely
\[
A_L\le \varepsilon_1 k^2. \tag{1}
\]
Next unfold the medium event
\(\noderef{TwoBiteMediumClosedPairsEvent}(k,\varepsilon,\varepsilon_1,\omega)\).
It says that for every \(k\)-set \(J\),
\[
\noderef{ClosedPairsBound}
\left(|\noderef{TwoBiteClosedPairsInClass}(J,M_J)|,\varepsilon_1,k\right)
\]
holds.  Applying it to the present set \(I\) and using the hypothesis
\(|I|=k\) gives
\[
\noderef{ClosedPairsBound}(A_M,\varepsilon_1,k),
\]
which, after unfolding \(\noderef{ClosedPairsBound}\), is
\[
A_M\le \varepsilon_1 k^2. \tag{2}
\]
The small event is used in the same way.  Unfolding
\(\noderef{TwoBiteSmallClosedPairsEvent}(k,\varepsilon,\varepsilon_1,\omega)\)
and applying it to \(I\) with \(|I|=k\) gives
\[
\noderef{ClosedPairsBound}(A_S,\varepsilon_1,k),
\]
hence
\[
A_S\le \varepsilon_1 k^2. \tag{3}
\]
Finally, the same-color huge hypothesis in the statement is already exactly
\[
\noderef{ClosedPairsBound}(A_H^{=},\varepsilon_1,k),
\]
with the same let-bound \(\mathsf{hugeRed}\) and \(\mathsf{hugeBlue}\)
predicates and the same two projected-pair sums as the conclusion.  Unfolding
\(\noderef{ClosedPairsBound}\) gives
\[
A_H^{=}\le \varepsilon_1 k^2. \tag{4}
\]

The real expression in the conclusion is exactly the four-term sum
\[
A_L+A_M+A_S+A_H^{=},
\]
with the terms ordered as large, medium, small, and then the parenthesized
huge same-color projected contribution.  Applying monotonicity of addition to
(1), (2), (3), and (4) in that order gives the line-by-line chain
\[
\begin{aligned}
A_L+A_M+A_S+A_H^{=}
&\le
(\varepsilon_1k^2)+(\varepsilon_1k^2)+A_S+A_H^{=}\\
&\le
(\varepsilon_1k^2)+(\varepsilon_1k^2)+(\varepsilon_1k^2)+A_H^{=}\\
&\le
(\varepsilon_1k^2)+(\varepsilon_1k^2)+(\varepsilon_1k^2)
  +(\varepsilon_1k^2)\\
&=4\varepsilon_1k^2.
\end{aligned}
\]
Thus
\[
A_L+A_M+A_S+A_H^{=}
\le 4\varepsilon_1k^2.
\]
Since \((4\varepsilon_1)k^2=4\varepsilon_1k^2\) by associativity of real
multiplication, this is precisely
\(\noderef{ClosedPairsBound}\) for the displayed Lean sum with parameter
\(4\varepsilon_1\) and scale \(k\).
\end{proof}
